% ECONOMICS_PROBLEM: {"id": "C73-5", "title": "Undiscounted equilibrium in three-player recursive games", "jel_code": "C73", "source_id": "OP-0157"}
% !TeX program = xelatex
\documentclass[11pt,a4paper]{article}
\usepackage[margin=23mm]{geometry}
\usepackage{fontspec}
\setmainfont{Latin Modern Roman}
\usepackage{amsmath,amssymb,mathtools}
\usepackage{xcolor,enumitem,booktabs,longtable,array,xurl,hyperref}
\definecolor{muted}{HTML}{53636C}
\hypersetup{colorlinks=true,urlcolor=blue,linkcolor=blue}
\setlength{\parindent}{0pt}
\setlength{\parskip}{5pt}
\newcommand{\R}{\mathbb R}
\newcommand{\E}{\mathbb E}
\newcommand{\N}{\mathbb N}
\newcommand{\ind}{\mathbf 1}
\newcommand{\Var}{\operatorname{Var}}
\newcommand{\Cov}{\operatorname{Cov}}
\newcommand{\supp}{\operatorname{supp}}
\DeclareMathOperator*{\argmin}{arg\,min}
\DeclareMathOperator*{\argmax}{arg\,max}
\newcommand{\entrymeta}[3]{{\sffamily\small\color{muted}\textbf{\texttt{#1}}\quad|\quad #2\par #3\par}}
\newcommand{\Problemref}[1]{\texttt{#1}}
\newcommand{\JELref}[2]{\texttt{#2} (JEL \texttt{#1})}
\begin{document}
\section*{Undiscounted equilibrium in three-player recursive games}
\textbf{Problem ID:} \texttt{C73-5}\quad \textbf{JEL:} \texttt{C73}\par
% BEGIN ECONOMICS STATEMENT
\label{problem:OP-0157}
% GAME_THEORY_PROBLEM: {"local_id": "GT-027", "original_number": 27, "kind": "P", "entry_kind": "statement_only", "published": false, "review_required": true, "category": "game_theory"}
\label{game:GT-027}
{\sffamily\small\color{muted}
{\sffamily\small\color{muted}\textbf{GT-027} | Stochastic, repeated, and dynamic games\par P | Historical source-reported existence question; later status not exhaustively audited\par}
\par}
\subsection*{Definitions and mathematical statement}
Let $N=\{1,2,3\}$. A finite recursive stochastic game has a finite state set $S$, a subset $B\subseteq S$ of absorbing states, finite nonempty action sets, and a transition kernel. Every $b\in B$ remains at $b$ under all actions and has a fixed bounded reward vector $g(b)\in\mathbb R^3$. At every nonabsorbing state every stage reward is zero. States and joint actions are observed; strategies may randomize and depend on the entire finite history.

Let $\theta=\inf\{t\geq0:s_t\in B\}$, with $\inf\varnothing=\infty$, and define
\[
 \Gamma_i(\sigma)=\mathbb E_{s_0,\sigma}
 \left[g_i(s_\theta)\mathbf1\{\theta<\infty\}\right],
\]
where the random variable is zero on $\{\theta=\infty\}$. The question is
\[
 \forall G\ \forall s_0\ \forall\varepsilon>0\quad
 \exists\sigma\quad\forall i\in\{1,2,3\}\ \forall\tau_i,\qquad
 \Gamma_i(\tau_i,\sigma_{-i})\leq\Gamma_i(\sigma)+\varepsilon.
\]
Here the pathwise Ces\`aro average converges to the terminal random payoff (or zero if absorption never occurs); bounded convergence identifies its expected limit with $\Gamma_i$.
\subsection*{Short English statement}
Must every finite three-player recursive game have an approximate equilibrium for its terminal absorption payoff?
\subsection*{Variants and scope boundaries}
Negative absorbing rewards are allowed. An absorbing game with only one nonabsorbing state is a proper subclass; a positive recursive game is another restriction. Uniform finite-horizon equilibrium is a separate requirement. Strategies are unrestricted behavioral strategies: absence of a stationary equilibrium does not establish absence of an unrestricted approximate equilibrium.
\subsection*{Source relationship and status}
[S1], the introduction on p. 83, explicitly singles out three-player recursive-game approximate existence as open. Its stationary counterexamples concern a different strategy restriction. The 2019 statement is preserved with an honest source-date limitation.
\subsection*{References}
\begingroup\footnotesize\raggedright
\textbf{[S1]} Kristoffer Arnsfelt Hansen and Mikhail Raskin, \emph{A Stay-in-a-Set Game without a Stationary Equilibrium}, EPTCS 305 (2019), pp. 83--90, doi:10.4204/EPTCS.305.6; arXiv:1903.11935. Locator: introduction, p. 83, paragraph on undiscounted approximate equilibrium in recursive games. \url{https://arxiv.org/pdf/1903.11935}
\par\endgroup

\subsection*{Common conventions: Source mathematical conventions}

\textbf{Finite games.} For a finite set $A$, $\Delta(A)$ is its probability
simplex. Players choose independent mixed strategies in Nash equilibrium;
correlation is allowed only when explicitly part of the solution concept.
In a game with payoff functions $u_i$ and strategy profile $x$, an additive
$\varepsilon$ Nash equilibrium satisfies
\[
 u_i(x)\geq u_i(x_i',x_{-i})-\varepsilon
 \quad\text{for every player $i$ and every admissible deviation $x_i'$}.
\]
For numerical approximation, payoffs are normalized as locally specified.
An $\varepsilon$ payoff residual is distinct from distance at most
$\varepsilon$ to the set of exact equilibria. Point convergence, convergence
of empirical play, and convergence in distance to an equilibrium set are
also distinct.

\textbf{Computational representations.} Unless stated otherwise, a complexity
question takes rational data with binary encoding; polynomial time is measured
in total bit length. A fixed-accuracy polynomial algorithm is not necessarily
polynomial in $1/\varepsilon$, and neither is necessarily polynomial in
$\log(1/\varepsilon)$. Succinct games and value, demand, or sampling oracles
must declare their interfaces. Exact real solutions require an explicit
representation; an unspecified list of arbitrary real numbers is not a
finite computational output. A claimed algorithmic barrier is meaningful
only for its specified model and reductions.

\textbf{Dynamic games.} Histories, information, observation, and allowable
randomization are part of the model. A stationary strategy depends only on
the current state; a positional strategy in a deterministic graph game
chooses one action at each controlled vertex. Nash equilibrium at an initial
state and equilibrium after every history are different requirements.
An expectation of a pathwise $\liminf$ payoff, a $\liminf$ of expected averages,
a limiting expected average, and a uniform finite-horizon equilibrium are
different criteria. None is silently substituted for another.

\textbf{Allocation and mechanisms.} A complete allocation assigns every item
exactly once unless otherwise stated. Zero-valued goods, negative values,
capacity constraints, feasibility, lotteries, and copies of goods affect
fairness definitions and are specified locally. Dominant-strategy incentive
compatibility, Bayesian incentive compatibility, universal truthfulness,
and truthfulness in expectation impose different inequalities. Whenever
an objective ratio has zero denominator, its convention is stated or those
instances are excluded.

\textbf{Infinite-dimensional models.} Expectations, integrals, stochastic
equations, and optimization objectives require their local regularity,
integrability, filtrations, and admissible domains. A backward--forward
mean-field system is a boundary-value problem; it is not an initial-value
semiflow without an additional construction. Long-time, large-population,
and vanishing-noise limits require an order of limits and a convergence
topology. If a minimum need not be attained, the objective is its infimum
or supremum and the approximation requirement is stated separately.
% END ECONOMICS STATEMENT
\end{document}
